% =====================================================================
% Chapter 5, part 1: Discussion. Transcribed from the reviewed
% DISCUSSION.html % reviewed draft; pooled-only framing.
% =====================================================================

\section{Discussion}
\label{sec:discussion}

The results chapter measured what this dataset supports under
subject-independent validation: subject-level scores that sit in the same
band as the score recording duration provides on its own
(Section~\ref{sec:res-floors}). Two questions follow from that measurement,
and this chapter answers both with experiments rather than argument. The
first is why published radar gait studies report so much more.
Section~\ref{sec:disc-protocol} reproduces the literature's numbers on this
very dataset by changing nothing but the validation design, first with the
campaign's own network and then with the field's favourite architecture, and
shows that those numbers collapse the moment a test subject is genuinely
new. The second is what the scores obtained here are actually made of.
Section~\ref{sec:disc-attention} reads the models' own attention maps,
Section~\ref{sec:disc-intervention} deletes what they attend to and confirms
that no hidden gait signature was underneath, and
Section~\ref{sec:disc-positive} states the findings the data does establish.
Section~\ref{sec:limitations} closes with the limitations; the concrete path
to stronger results follows in the future-work section.

The framing this chapter defends can be stated in advance. The gap between
this thesis's numbers and the literature's is not a performance gap; it is a
\textbf{measurement gap}, because the two are answering different questions.
The literature's question, whether a model can label recordings of people it
has already seen, is answered here too, and with literature-grade numbers.
The clinical question, whether a model can flag a person it has never seen,
is the one this thesis insists on, and its honest answer on this dataset is:
not yet, and not demonstrably from gait shape.

% ---------------------------------------------------------------------
\subsection{The random-split demonstration: the gap is the protocol}
\label{sec:disc-protocol}

The claim that the literature's numbers are inflated by validation design
would normally rest on reading papers. Here it was tested directly: the same
network and data as the main campaign, trained the field's way, with a
random hold-out split, early stopping on window-level scores and no window
weighting, so that the whole published recipe is reproduced rather than
argued about. Table~\ref{tab:protocol_ladder} reports the outcome.

\begin{table}[H]
  \centering
  \caption[The validation ladder]{One architecture under four validation
  designs. All rows use the ResNet-18 probe of
  Section~\ref{sec:deep}; the first three are trained identically and differ
  only in what the split keeps apart.}
  \label{tab:protocol_ladder}
  \small
  \begin{tabular}{@{}L{5.1cm}cL{7.2cm}@{}}
    \toprule
    Validation design & AUC & What it measures \\
    \midrule
    Window-level random split & 0.774 &
      Nearly identical overlapping windows of the same walk sit on both
      sides of the split (3 repeats). \\
    Recording-level random split (the literature's hold-out) & 0.760 &
      Every test recording comes from a person the model trained on
      (5 repeats). \\
    Same trained model, subjects excluded from training & 0.579 &
      The person is unknown; only the disease could help (3 repeats). \\
    Leave-one-subject-out (the thesis protocol) & 0.654 &
      Every subject scored by a model that never saw them
      (subject-level AUC, Table~\ref{tab:deepauc}). \\
    \bottomrule
  \end{tabular}
\end{table}

The rows deserve unpacking, because the mechanics decide what each number
means. In the \textbf{window split}, the three-second windows themselves are
shuffled into training, validation and test piles. In the \textbf{recording
split}, whole walks are shuffled instead, so each test walk is new but the
walker's other walks are in training; this is the design most published
studies use. In the \textbf{third row}, eight subjects (five control, three
PD) are set aside before anything is trained; the model trains
recording-split style on the remaining fifty subjects, and the finished
model then classifies every window of the eight strangers' recordings, with
a fresh eight in each repeat. In the \textbf{last row}, one subject is set
aside, the model trains on the other 57, and scores every window of all of
that person's recordings; the average of those window probabilities becomes
that one person's score, and repeating this 58 times gives the subject-level
AUC of the main campaign. One unit caveat applies: the first three rows are
window-level AUC averaged over the repeats, the level at which the
literature scores, while the last row is subject-level, one score per
person. The trend is the point: the more the split leaks, the higher the
number, and stripping the leakage returns the same model, on the same data,
to the modest band reported in Chapter~\ref{ch:results}.

\begin{figure}[H]
  \centering
  \includegraphics[width=\textwidth]{disc_collapse.png}
  \caption[The collapse test]{The collapse test, the third row of
  Table~\ref{tab:protocol_ladder}: three instances of the ResNet-18 probe, one
  line each. The left end of each line is that model's window AUC on test
  recordings of subjects it saw in training; the right end is the same model
  on the eight subjects excluded from training.}
  \label{fig:disc_collapse}
\end{figure}

Each line in Figure~\ref{fig:disc_collapse} is one trained model, not one
subject. On the left, every model scores high on test recordings of people
whose other walks it trained on; on the right, the very same models are
asked about the eight people they have never seen, and every line falls,
with two of the three landing at chance. A model that had learned the
disease would hold its score on strangers, because Parkinson's looks like
Parkinson's in people it never met. A model that learned \textbf{people}
scores high on the familiar and collapses on the unknown, and that is what
happens. The inflated numbers are real numbers; they are simply answers to
the recognition question, not the screening question. Once it is established
that the validation design alone manufactures a difference on the order of
0.2 AUC on this very dataset (0.18 to 0.20 between the averaged rows of
Table~\ref{tab:protocol_ladder}, and up to 0.4 in the most extreme repeat of
Figure~\ref{fig:disc_collapse}), comparisons with published accuracies stop
being embarrassing and start being diagnostic.

% ---------------------------------------------------------------------
\subsection{The same demonstration with the field's own network}
\label{sec:disc-alexnet}

The demonstration above uses the thesis's own network, and a sceptic could
object that the collapse says something about that particular model rather
than about the protocol. The experiment was therefore repeated with
\textbf{AlexNet}, the architecture behind the field's headline micro-Doppler
result \cite{hayashi2021}, given the same input adaptation as the ResNet-18
probe of Section~\ref{sec:deep}: ImageNet-pretrained, with the first
convolution averaged to the two radar channels. Unlike the probe, which
retrains a single 1{,}026-parameter layer, AlexNet's full classifier block is
retrained (54.5 million trainable parameters) at the source literature's
learning rate. The window
cache and augmentation are those of the main campaign; the random-split runs
additionally adopt the literature's training conventions, early stopping on
window-level scores and no window weighting. Table~\ref{tab:alexnet}
reports the outcome under both designs.

\begin{table}[H]
  \centering
  \caption[AlexNet under the two protocols]{AlexNet under the literature's
  protocol and under the thesis protocol. A seed is one repeat of the split:
  the same recordings shuffled with a different random starting point.}
  \label{tab:alexnet}
  \footnotesize
  \begin{tabular}{@{}llccc@{}}
    \toprule
    Validation design & Train / val / test & Subject AUC & 95\,\% CI &
      Seen in training \\
    \midrule
    Random split, seed 1 & 70/15/15\,\% recordings & 0.875 & -- & 100\,\% \\
    Random split, seed 2 & 70/15/15\,\% recordings & 0.891 & -- & 100\,\% \\
    Random split, seed 3 & 70/15/15\,\% recordings & 0.735 & -- & 100\,\% \\
    Leave-one-subject-out & 49/8/1 subjects, $\times$58 & 0.695 &
      [0.554, 0.833] & 0\,\% \\
    \bottomrule
  \end{tabular}
\end{table}

\begin{figure}[H]
  \centering
  \includegraphics[width=\textwidth]{disc_alex_ladder.png}
  \caption[AlexNet: one architecture, two protocols]{AlexNet's subject-level
  AUC under the literature's random split (red dots, one per seed, each
  labelled with its score) and under leave-one-subject-out (blue dot with
  its 95\,\% confidence interval), against the chance and duration-alone
  lines.}
  \label{fig:alexladder}
\end{figure}

Figure~\ref{fig:alexladder} admits three readings. First, the field's own
architecture reaches 0.87 and 0.89 the moment every test subject is someone
it trained on, and drops the moment they are not: \textbf{the split, not the
architecture, decides} whether the number is literature-grade or honest.
Second, the inflated number is itself unstable: reshuffling the same
recordings moves it from 0.89 down to 0.74, a swing of 0.16 between seeds.
Third, the honest interval reaches down past the duration floor of
Table~\ref{tab:floors}, so even the best protocol-compliant score in this
project cannot be distinguished from what elapsed time alone provides.

An AUC summarises how well the scores \emph{rank} people; a screening
decision is a yes or no. Figure~\ref{fig:alexconfusion} therefore shows the
same network's actual decisions twice, first on subjects it saw in training
and then on subjects it has never seen. The two panels count different
things: on the left, each cell counts three-second windows, every test
window classified on its own at the default 0.5 cut-off and pooled over the
three seeds' test sets; on the right, each cell counts people, with all
windows of a person's recordings scored, their average becoming that
person's single probability, and the person called PD when it exceeds the
prevalence threshold of 0.431 (Section~\ref{sec:res-decisions}).

\begin{figure}[H]
  \centering
  \includegraphics[width=\textwidth]{disc_alex_confusion_pair.png}
  \caption[AlexNet's decisions, seen versus never-seen]{AlexNet's decisions
  on subjects seen in training (left: random split, window level, pooled
  over the three seeds) and on never-seen subjects (right:
  leave-one-subject-out, one decision per person, threshold at the PD
  prevalence).}
  \label{fig:alexconfusion}
\end{figure}

On seen subjects the model is balanced and looks useful, calling about
70\,\% of each class correctly. On never-seen subjects the balance breaks:
it catches 21 of the 25 PD patients, but only by flagging 18 of the 33
healthy controls along with them, a specificity of 45\,\%, worse than a coin
flip. The left panel is what the literature's validation rewards, and it is
real; it is also the answer to the recognition question, because every one
of those windows belongs to a person whose other walks were in training. The
right panel is the clinical question: faced with a stranger, the network can
only lean toward calling recordings PD, so of the 39 people it would send
onward for review, 18 are healthy. The same architecture is clinic-ready on
people it knows and close to a biased guess on the population a screening
tool actually exists for.

Two details make the honest AlexNet score itself instructive. Its subject
scores carry \textbf{more duration residue than any of the campaign's
headline models} (correlation $+0.25$ with window count, against $+0.08$ for
the ResNet-18 probe): with 54.5 million parameters free, it drinks
more deeply from the same confound; it does not find a new signal. And its
best validation epoch averages 0.9, meaning the score is essentially in
place after a single pass over the data, the learning profile of an easy
shortcut rather than of a subtle clinical pattern.

% ---------------------------------------------------------------------
\subsection{What the models actually read}
\label{sec:disc-attention}

The results chapter reported the scores; the attention analysis explains
them. Grad-CAM maps \cite{selvaraju2017} were computed for every window of
each of the 58 held-out subjects, for both image models: the compact CNN
trained from scratch (SmallCNN) and the pretrained deep network (the
ResNet-18 probe). One gallery per model follows, so the finding can
be checked on each.

\begin{figure}[H]
  \centering
  \includegraphics[width=\textwidth]{disc_gradcam_smallcnn.png}
  \caption[Grad-CAM gallery, SmallCNN]{Where the SmallCNN's PD output looks,
  for confident correct decisions (top) and confident mistakes (bottom):
  spectrogram in grey, attention in colour.}
  \label{fig:gradcam_smallcnn}
\end{figure}

In every window of Figure~\ref{fig:gradcam_smallcnn} called PD, right or
wrong, the attention floods the empty background; the gait envelope itself
is the one region left unhighlighted. In every window called control the map
stays dark. The decision variable is a property of the recording, its
background and enhancement texture, not of the walk. The mistakes obey the
same rule, which is why they are systematic rather than random: noisy-backed
controls get flagged, and clean-backed patients get missed, exactly the
pattern already visible in the case gallery of Figure~\ref{fig:casegallery}.

\begin{figure}[H]
  \centering
  \includegraphics[width=\textwidth]{disc_gradcam_resnet.png}
  \caption[Grad-CAM gallery, ResNet-18 probe]{The same gallery for
  the ResNet-18 probe: confident correct decisions on top,
  confident mistakes below.}
  \label{fig:gradcam_resnet}
\end{figure}

The deep network in Figure~\ref{fig:gradcam_resnet} obeys the same rule,
expressed more diffusely: windows it calls PD are flooded with attention
across the whole image, silent rows included, while windows it calls control
are left almost unattended. Averaged over every window of all 58 held-out
subjects, its attention density on called-PD windows is 0.37 against 0.13 on
called-control windows; for the SmallCNN the same pair is 0.33 against 0.04.
Two different architectures, one trained from scratch and one built on
frozen ImageNet features, converge on the same behaviour: respond to the
overall character of the recording, not to any localised gait structure.

Three quantitative facts, computed over every window of all 58 held-out
subjects, complete the picture:

\begin{itemize}
  \item For the SmallCNN, attention density is highest in the silent rows
        above 500~Hz (0.19 per pixel against 0.12 in the torso band), a
        region containing no gait energy, and its attention sits visibly
        higher in frequency on PD windows than on control windows (centre of
        mass 380 against 333~Hz).
  \item For the SmallCNN, turn-straddling windows draw \textbf{four to six
        times} the attention of steady-walking windows, matching the
        ablation of Table~\ref{tab:ablations} in which excluding the turn
        collapsed its score.
  \item The probe differs in one telling way: because its backbone
        is frozen, its maps average out \textbf{identical between the true
        classes} (background attention share 0.337 on PD windows against
        0.338 on control windows). A fixed ImageNet description of these
        images contains nothing class-specific to point at, so its thin
        trained head can only read overall texture.
\end{itemize}

Taken together, the mechanism of every score in this thesis is now
explicit: \textbf{elapsed time}, carried by turn-and-stand content and
window counts, plus \textbf{per-recording character}, carried by background
and enhancement texture. Both are properties a screening tool must not rely
on; neither is the spectral shape of a Parkinsonian walk.

% ---------------------------------------------------------------------
\subsection{The intervention test: removing what they read}
\label{sec:disc-intervention}

An interpretability finding invites an intervention: if the models read the
background, delete the background and see what remains. The follow-up
experiment, pre-registered before any window was rebuilt, subtracted each
recording's background profile and cropped the silent rows beyond 500~Hz,
then re-ran both image models under the unchanged protocol.
Table~\ref{tab:bgsuppressed} reports the outcome.

\begin{table}[H]
  \centering
  \caption[Background-suppressed results]{Both image models with the background suppressed (pooled AUC,
  full leave-one-subject-out protocol).}
  \label{tab:bgsuppressed}
  \small
  \begin{tabular}{@{}lccc@{}}
    \toprule
    Model & Standard input & Background-suppressed & 95\,\% CI (suppressed) \\
    \midrule
    SmallCNN        & 0.616 & 0.624 & [0.471, 0.768] \\
    ResNet-18 probe & 0.654 & 0.686 & [0.539, 0.824] \\
    \bottomrule
  \end{tabular}
\end{table}

Two conclusions follow, both informative:

\begin{itemize}
  \item \textbf{No hidden gait signal surfaced.} If the texture had been
        masking a real gait signature, removing it should have lifted the
        scores clearly above the duration floor. Every confidence interval
        still contains 0.610.
  \item \textbf{The nuisance is distributed, not localised.} The scores did
        not collapse either: enough per-recording character survives inside
        the gait band, plausibly in the per-recording contrast enhancement,
        for the models to keep finding it. The recording-level nuisance is a
        property of how each file was processed, not one removable region of
        the image.
\end{itemize}

% ---------------------------------------------------------------------
\subsection{What was found, stated positively}
\label{sec:disc-positive}

Not everything in this thesis is a negative result. The findings of
Table~\ref{tab:positive} are established by the data and hold under the
strict protocol; each row gives the finding, the evidence in brief, and what
it tells us.

Against the literature, these results are not an outlier; they are what the
few honest data points predict. Of the sixteen works reviewed in
Chapter~\ref{ch:soa}, only two report a genuinely subject-independent protocol, and one
of them saw its scores fall the moment it moved from a pooled split to
leave-one-subject-out evaluation \cite{fard2026}. The study closest in modality to
this work preferred its modest envelope model over a network trained on the full image
whose higher score its own authors distrusted, because that network had
keyed on background differences between recording sites \cite{hayashi2021}.
This thesis extends that honest line to Parkinson's screening and explains,
mechanistically, where the distrusted scores come from.

\begin{table}[H]
  \centering
  \caption[The positive findings]{The positive findings.}
  \label{tab:positive}
  \small
  \begin{tabular}{@{}L{3.6cm}L{5.1cm}L{5.4cm}@{}}
    \toprule
    \textbf{Finding} & \textbf{Evidence} & \textbf{What it tells us} \\
    \midrule
    \textbf{PD subjects take measurably longer}, and the extra time
    concentrates in standing up and turning. &
    The duration gap between the groups is clear in the recording lengths,
    is largest in the walk variant that includes the chair phase, and
    shrinks when the turn is removed. &
    Radar captures the slowness of movement that clinicians already test
    for; the timing of a walk is a real, usable signal. \\
    \addlinespace
    \textbf{Radar timing is a credible screening quantity.} &
    Recording duration alone separates the groups above chance under honest,
    subject-independent validation. &
    An automated timed-up-and-go measurement is this dataset's trustworthy
    first-line signal, and it needs no deep network at all. \\
    \addlinespace
    \textbf{A walk identifies the walker.} &
    Recordings of the same person are far more alike than recordings of
    different people, for every one of the 58 subjects
    (Section~\ref{sec:res_fingerprint}). &
    Subject-independent validation is mandatory, and radar gait could serve
    as a biometric even where diagnosis fails. \\
    \addlinespace
    \textbf{The field's validation gap is real and measurable.} &
    Relaxing the protocol on this one dataset reproduces literature-grade
    numbers that collapse on unseen subjects. &
    Published accuracies reflect validation design as much as clinical
    signal, so honest evaluation has to be demanded before numbers are
    compared. \\
    \bottomrule
  \end{tabular}
\end{table}

% =====================================================================
\section{Limitations}
\label{sec:limitations}

Four limitations bound what this thesis can claim.
Table~\ref{tab:limitations} states each one together with its consequence
and the item that would resolve it.

\begin{table}[H]
  \centering
  \caption[Limitations, and what would resolve each]{Limitations, and what would resolve each.}
  \label{tab:limitations}
  \small
  \begin{tabular}{@{}L{4.4cm}L{9.9cm}@{}}
    \toprule
    \textbf{Limitation} & \textbf{Consequence, and what would resolve it} \\
    \midrule
    \textbf{No age metadata.} &
    The dominant untestable confound: PD subjects are typically older, radar
    separates young from elderly walkers at 94.9\,\% \cite{hayashi2021}, and
    even the surviving duration signal could partly be ageing. Subject ages
    would allow an age-matched analysis. \\
    \addlinespace
    \textbf{Combined foot channels hide left/right asymmetry.} &
    Asymmetry is an established PD marker, and it is structurally invisible
    after the two foot radars are merged into one channel. The original
    per-sensor recordings would make it measurable. \\
    \addlinespace
    \textbf{58 subjects.} &
    With 58 people, every score carries a wide uncertainty range, so a small
    real effect can neither be proven nor ruled out at this size. Only a
    much larger group, in the hundreds, can settle small differences. \\
    \addlinespace
    \textbf{No clinical information beyond the label.} &
    Each subject is only marked PD or control. Without disease stage,
    symptom severity or medication state, it is impossible to test whether
    the signal grows with severity, or how treatment changes it, both of
    which matter for a screening tool. \\
    \bottomrule
  \end{tabular}
\end{table}
